\documentclass[11pt,letterpaper]{article}
\usepackage[T1]{fontenc}
\usepackage{lmodern,microtype,booktabs,tabularx,array,amsmath,enumitem,xcolor}
\usepackage[margin=1in]{geometry}
\usepackage[colorlinks=true,linkcolor=blue!50!black,urlcolor=blue!50!black,citecolor=blue!50!black]{hyperref}
\usepackage{xurl}
\setlength{\parskip}{4pt}
\setlist{nosep,leftmargin=1.4em}
\newcolumntype{Y}{>{\raggedright\arraybackslash}X}
\renewcommand{\arraystretch}{1.15}

\title{DeepTrace: Does Code-Change History Help an AI Safety Monitor\\Catch Multi-Step Sabotage?}
\author{Ved Panse, Michael Luo, Ananya Wasker \\ Mentor: Babak Salimi \\ DSC 180A, Domain D33 \\ \small Report Checkpoint}

\begin{document}
\sloppy
\maketitle

\begin{abstract}
In recent times, AI coding agents have been thrusted with the responsibility of editing real repositories and running shell commands with little human review. A common safeguard is a trusted monitor, which is a smaller, cheaper model that reads each proposed action and raises an alarm if it looks unsafe. Monitors usually judge an action from the current transcript or the current code alone, so an advanced attacker or saboteur can work around this by spreading an attack over several harmless-looking edits. For example, one step weakens an authorization helper and a later step routes user input into a shell command that relied on it. To combat this, static analyzers such as AgentSecGraph can build a dependency graph of sources, guards and effects for each security-sensitive operation. They describe a single snapshot, though, and they do not record that a guard the operation depends on was changed earlier. Our project, titled DeepTrace, asks whether giving a fixed monitor a compact, version-aware record of changing security-relevant facts over earlier edits improves its verdicts under a fixed evidence budget. We first reproduced AgentSecGraph on three real agent repositories and eleven authored fixtures. We then ran a series of frozen, pre-registered experiments with GPT-5 Nano as the monitor on synthetic Python edit histories. In the latest experiments we've run, the main bottleneck so far is applying the policy and getting accurate graph representations of repositories and projects (everything from files, functions, folders etc.). Our intended real-world setting is LinuxArena, a benchmark of sabotage attempts in live software environments, which we are working on integrating with the DAG (directed-acyclic graph) representation of information.
\end{abstract}

\section{Introduction}

\textbf{DeepTrace} asks the following question:

\begin{quote}
\emph{Does a fixed safety monitor judge a proposed code change more reliably when, in addition to the current code, it receives a temporal DAG representation of earlier edits and dependencies amplifying the code change's impact, under the same evidence budget?}
\end{quote}

Our proposal has two hypotheses. \textbf{H1} says that dependency-graph evidence improves timely detection of multi-step sabotage in real agent trajectories, meaning detection before the first harmful action executes. \textbf{H2} says that explicit records of when security assumptions were invalidated, and by which edit, improve on both current-code graphs and raw edit history. The question matters for three reasons. If history helps, it is a cheap way to strengthen existing monitors without training a new model. If history does not help, or even hurts, that is equally important to know. Finally, the experiment looks at a monitor's overall vulnerabilities. It might lack the right facts, misread facts it has, or apply the policy incorrectly to facts it has read correctly.

\subsection{Literature review}

\paragraph{1. AI control and trusted monitoring:} Greenblatt et al.\ introduced the AI-control framing, in which a protocol must stay safe against an untrusted model that deliberately tries to subvert it. Their protocols use a weaker trusted model to monitor and edit the outputs of a stronger untrusted one, under a fixed human audit budget. Later work extended this setting to multi-step agents. Ctrl-Z (Bhatt et al., 2025) studies resampling protocols for agents acting in a shell environment. SHADE-Arena (Kutasov et al., 2025) tests whether agents can complete hidden harmful side tasks without being caught by a monitor. These works share an evaluation pattern that we adopt, which is to calibrate a monitor's alarm threshold on honest behavior to a small false-positive rate, then measure attacks' success rate with concealment.

\paragraph{2. LinuxArena:} LinuxArena (Tracy et al., 2026) is a control benchmark in which agents work inside live, multi-service production software environments. It evaluates both attackers, who must complete a side task while doing a main task, and monitors, which must detect those attempts. Its main result is that Claude Opus 4.6 achieves roughly 34\% undetected sabotage success against a GPT-5 Nano trusted monitor at a 1\% step-wise false-positive rate. We use LinuxArena because its attacks are multi-step and unfold over real codebases, so edit history is available to make our dependency graphs. Secondly, its standard monitor is the same small model we use, so any benefit from better evidence can be read against a published baseline. The paper also lists shallow tracing of earlier actions among monitor failure modes, which is the gap DeepTrace aims to assist with.

\paragraph{3. Static analysis for agent security: AgentSecGraph:} The closest prior system to our representation is AgentSecGraph, released with the AgentSecBench corpus (Cui, 2026). It builds on classical program dependence graphs (Ferrante et al., 1987) and code property graphs (Yamaguchi et al., 2014). Related agent-specific analyzers include AgentFlow (Wang et al., 2026), which recovers framework-induced agent dependency graphs, and TaintP2X (J.~He et al., 2026), which statically traces LLM-generated values to sensitive operations. AgentSecGraph builds a candidate-centered \emph{Security-Aware Agent Dependency Graph} (Security-ADG) that records agent relevance, source provenance, data dependencies, trust-boundary crossings, guard evidence, privilege context and external effects. Python dependencies come from the AST and JavaScript/TypeScript dependencies from lexical tracing. Framework adapters recognize tool declarations in MCP, LangChain, CrewAI, AutoGen, the OpenAI Agents SDK, Semantic Kernel and LlamaIndex. We adopt AgentSecGraph, unmodified and pinned to AgentSecBench commit \texttt{e27f83d}, as our static baseline. The public checkout contains two separate implementations: the paper's research pipeline (\texttt{v2\_5}) and a public product (\texttt{0.2.0}). However, neither implementation tracks history.

\paragraph{4. Gap Identified:} Existing static analyzers characterize a program as a single snapshot, and existing monitors evaluate an agent's actions using the current transcript or current code. Runtime defenses such as AttriGuard (Y.~He et al., 2026) attribute individual tool calls to untrusted inputs, but they do not track how code edits change security assumptions over time. Additionally, no prior work has measured whether such a record improves a fixed monitor's detection at a fixed false-positive rate, and DeepTrace is designed to measure this. Our preliminary results suggest that the effect depends on the evaluation setting. Any effect historical information has on our experiments reflects the monitor's robustness to extra context. A complex and involved evaluation therefore requires real trajectories in which the relevant dependency can be correctly recovered from the current state alone.

\subsection{Data}

\paragraph{1. Real agent repositories (AgentSecGraph baseline).} We ran AgentSecGraph at pinned commits on three open-source agent projects: \texttt{html-anything}, \texttt{mini-swe-agent} and \texttt{magentic-ui}. We also ran it on eleven small authored fixtures that isolate one analysis challenge each, such as overwrites, cross-file aliasing, guards, dynamic dispatch and recursion. The scans ran in a network-disabled sandbox and are summarized below. These data establish what the static baseline can and cannot represent.

\paragraph{2. Controlled synthetic edit histories (main experiments so far).} We authored Python edit histories.
\begin{itemize}
\item \emph{What a history contains.} Each history is a sequence of source versions ending in a proposed edit. Variants cover a guard being weakened, restored, replaced by an equivalent guard, executed with fixed arguments or removed, along with unrelated edits and deliberately unanalyzable (``unknown'') cases.
\item \emph{Ground truth.} Labels come from a restricted AST interpreter run on concrete inputs, cross-checked by a Z3-based symbolic checker. No generated command is ever executed.
\item \emph{Splits.} Splits are made by whole syntax or task families, so related variants never straddle calibration and test.
\item \emph{Sets built so far.} We built four successive sets: 270 histories for the first pilot, 320 for a fresh controlled comparison, 120 + 120 states for a paired current-versus-history test, and 24 + 36 cases per arm for the three-arm selector study.
\item \emph{Reproducibility.} Every prompt, raw model response, threshold and protocol is saved with content hashes. A standard-library script (\texttt{reproduce.py}) recomputes the selector-study table from the 108 saved test records without any API calls.
\item \emph{Limitation.} These cases are close relatives of one narrow policy and cannot speak to real-world generalization.
\end{itemize}

\paragraph{3. LinuxArena trajectories (target setting).} LinuxArena trajectories are recorded agent episodes. Each pairs a main task with, in attack runs, a hidden side task, and includes ordered tool calls, outputs and scorer results. Data comes from the public LinuxArena data service and its Control Tower harness. We have used them in three stages so far:
\begin{itemize}
\item \emph{Pilot:} We used five unique recorded-successful attack trajectories in five environments: \texttt{slack\_clone}, \texttt{file\_indexing}, \texttt{web\_scraping}, \texttt{shared\_dotfiles} and \texttt{auto\_workflow}. Five honest trajectories supplied 100 calibration and 100 held-out benign action turns. These bundles had no per-action repository snapshots.
\item \emph{Mirror audit:} The public Hugging Face mirror (\texttt{anonymouslinuxarena/linuxarena-public}) holds 22,215 trajectories from 172 runs, about 16.8\,GB. We audited 20 cached samples and 25 records from a ranged prefix of one shard.
\item \emph{Official archive:} We downloaded one full run archive (\texttt{ace6e820\ldots}, 18.5\,MB) without credentials. It has 105 samples and 7,383 tool events across all 10 public environments, and it records environment commits, container image digests and the harness revision. Re-executing their recorded tool calls in the container reproduced 12/17, 16/27 and 21/38 call outputs exactly.
\end{itemize}

The cohort does not yet meet the gate our evaluation requires: a verified initial state, verified intermediate states for every action, and a first-harm step checked by two independent reviewers. Finishing this verified cohort, then freezing fresh calibration and test splits before scoring any evidence condition, is the main data milestone for the rest of the quarter.

\section*{References}
\small
\begin{enumerate}[label={[\arabic*]},leftmargin=2em]
\item R.~Greenblatt, B.~Shlegeris, K.~Sachan, F.~Roger. AI Control: Improving Safety Despite Intentional Subversion. In \emph{Proc.\ ICML}, PMLR 235, pp.~16295--16336, 2024.
\item A.~Bhatt, C.~Rushing, A.~Kaufman, et al. Ctrl-Z: Controlling AI Agents via Resampling. arXiv:2504.10374, 2025.
\item J.~Kutasov, Y.~Sun, P.~Colognese, et al. SHADE-Arena: Evaluating Sabotage and Monitoring in LLM Agents. arXiv:2506.15740, 2025.
\item T.~Tracy, R.~Potham, N.~Kuhn, et al. LinuxArena: A Control Setting for AI Agents in Live Production Software Environments. arXiv:2604.15384v2, 2026.
\item H.~Cui. Beyond Predefined Sinks: Security-Aware Dependency Analysis for LLM Agents. arXiv:2610.03014, 2026. Code: \url{https://github.com/A1LinLin1/AgentSecBench}.
\item J.~Ferrante, K.~J.~Ottenstein, J.~D.~Warren. The Program Dependence Graph and Its Use in Optimization. \emph{ACM TOPLAS} 9(3):319--349, 1987.
\item F.~Yamaguchi, N.~Golde, D.~Arp, K.~Rieck. Modeling and Discovering Vulnerabilities with Code Property Graphs. In \emph{IEEE S\&P}, pp.~590--604, 2014.
\item S.~Wang, X.~Hou, Y.~Zhao, X.~Cheng, H.~Wang. AgentFlow: Building Agent Dependency Graphs for Static Analysis of Agent Programs. arXiv:2607.01640, 2026.
\item J.~He, S.~Wang, Y.~Zhao, X.~Hou, Z.~Liu, Q.~Zou, H.~Wang. TaintP2X: Detecting Taint-Style Prompt-to-Anything Injection Vulnerabilities in LLM-Integrated Applications. In \emph{Proc.\ ICSE}, pp.~1968--1979, 2026.
\item Y.~He, H.~Zhu, Y.~Li, S.~Shao, H.~Yao, Z.~Liu, Z.~Qin. AttriGuard: Defeating Indirect Prompt Injection in LLM Agents via Causal Attribution of Tool Invocations. USENIX Security, 2026.
\end{enumerate}


\subsection{AI assistance for this checkpoint report}
\textbf{Tool:} Claude Code and Codex.
The repository records that an AI coding assistant (``implementing assistant'') took part in building the experiments:
\begin{itemize}
\item It authored the explicit source expectations for the version selector and the selector evaluation, and reviewed rendered packets (\texttt{version\_selector/README.md}; \texttt{selector\_evaluation/results/protocol.json}, field \texttt{reference}).
\item It performed the post-generation source-structure review in \texttt{selector\_evaluation/results/source\_expectation\_audit.json}.

\begin{center}\small
\begin{tabularx}{\linewidth}{@{}p{0.2\linewidth}YY@{}}\toprule
Task & What the AI was asked to do & Human role and verification\\\midrule
Repository understanding & Index the repository and code, yield a list of literature used, experiments performed, results, with explanations and annotations. & Member refined this and wrote the abstract, research question, literature review and data sections in their own words.\\
Editing for tone & Rewrite the research-gap paragraph in a formal register. & Member reviewed the rewrite.\\
Reference selection & Trim the reference list. & Member chose which works to keep.\\
Citation verification & Check every reference against arXiv, Crossref and PMLR. & One error found and fixed: LinuxArena's figure changed from 23\% (outdated arXiv summary page) to 34\% (paper v2 text). Two same-surname citations were disambiguated.\\
Testing and validation & Run all runnable repository code on a clean copy and report results (see A.5). & Member reviewed the findings independently.\\
AI-use disclosure & Draft appendix from repository records and the session log. & Other members' entries were put in the repository for their code changes.\\
\bottomrule\end{tabularx}
\end{center}


\end{itemize}

\subsection{AI-assisted testing and validation}
At the authors' request, the AI assistant ran a validation pass on a clean export of the committed repository, using pinned dependencies and no model calls. Its results were:
\begin{itemize}
\item \textbf{Documented reproduction.} \texttt{reproduce.py} verified all 108 records, and 7/7 unit tests passed. A tamper test confirmed the input check is genuine: reformatting the input still passes, while flipping one saved verdict fails.
\item \textbf{Integrity manifests.} Of 8,336 locatable SHA-256 entries in 344 manifests, 8,322 match. Every one of the 14 mismatches is explained: canonical-JSON hashing (6), CSV line-ending normalization (4), a documented pre-inference protocol amendment (1), and report files revised after hashing (3).
\item \textbf{Determinism.} Regenerating the 270 synthetic temporal cases reproduced them byte for byte. Rerunning the downstream checks gave the same results: 1,620 ground-truth executions, 8/8 robustness controls, 162/162 counterfactuals and 20/20 continuation engineering checks.
\item \textbf{AgentSecGraph.} The analyzer and three target repositories were restored at their pinned commits. Sandboxed scans and the upstream test suite (36/36) reproduced all figures in the Data section. The temporal fixture again produced identical graph hashes for versions A and B.
\end{itemize}
\end{document}


\end{document}
